\section{Eliminating every square relation}

The local row becomes global because the collision parameters have the
full permutation group.  We record the elementary binary linear algebra so
that the rank conclusion does not rest on a black-box module computation.

\begin{lemma}[Pair annihilator]\label{lem:pair-annihilator}
Let \(W\subset\F_2^9\) be the span of the vectors \(e_i+e_j\) with
\(i\ne j\).  Then \(W\) is the even-weight subspace, \(\dim W=8\), and
\[
 W^\perp=\langle\ones\rangle,
 \qquad \ones=(1,\ldots,1).
\]
\end{lemma}

\begin{proof}
Every pair vector has even weight, so \(W\) lies in the codimension-one
even-weight subspace.  The eight vectors \(e_1+e_j\), \(2\le j\le9\),
are independent, proving equality and \(\dim W=8\).  A vector orthogonal
to every \(e_i+e_j\) satisfies \(r_i+r_j=0\) for all pairs, hence all its
coordinates agree.
\end{proof}

\begin{proposition}[Only one candidate relation]\label{prop:one-relation}
For the square-relation kernel \(R\) in \eqref{eq:R},
\[
 R\subseteq\{0,\ones\}.
\]
\end{proposition}

\begin{proof}
If \(r\in R\), then the valuation of
\(\prod_i\beta_i^{r_i}\) is even at every discrete valuation of \(L\).
Applying Proposition~\ref{prop:odd-hill} gives
\(r_1+r_2=0\).  The subspace \(R\) is invariant under
\(\Gal(L/\Q)=S_9\).  Conjugating both the valuation and its labels supplies
\(r_i+r_j=0\) for every unordered pair.  Equivalently, \(R\subseteq
W^\perp\), and Lemma~\ref{lem:pair-annihilator} finishes the proof.
\end{proof}

To exclude \(\ones\), we need one standard fact about symmetric splitting
fields.

\begin{lemma}[Rational squares in an \(S_n\)-field]\label{lem:sign-field}
Let \(L/\Q\) be the splitting field of a separable polynomial \(f\) with
\(\Gal(L/\Q)=S_n\), \(n\ge3\).  For \(c\in\Q^\times\), if
\(c\in L^{\times2}\), then either \(c\in\Q^{\times2}\), or
\(\class{c}=\class{\Disc f}\) in \(\Q^\times/\Q^{\times2}\).
\end{lemma}

\begin{proof}
If \(c\) is not a rational square but becomes a square in \(L\), then
\(\Q(\sqrt c)\) is a quadratic subfield of \(L\).  Quadratic subfields
correspond to normal index-two subgroups of \(S_n\).  The unique such
subgroup is \(A_n\), because the abelianization of \(S_n\) is \(C_2\).
Its fixed field is the sign field, generated by the alternating
Vandermonde, and its square class is \(\Disc f\).  This argument depends
only on the discriminant square class and therefore also applies to a
nonmonic defining polynomial.
\end{proof}

The exact discriminant factorization is
\begin{equation}
D:=\Disc(P_9)
=2^{96}3^{12}13^3 19^5 41^3 59^5 9056471^2.                \label{eq:disc}
\end{equation}
Combining \eqref{eq:norm} and \eqref{eq:disc} gives
\begin{align}
 \class{N}_{\Q^\times/\Q^{\times2}}
   &=\class{3\cdot13\cdot19\cdot41\cdot59},              \label{eq:nclass}\\
 \class{D}_{\Q^\times/\Q^{\times2}}
   &=\class{13\cdot19\cdot41\cdot59},                    \label{eq:dclass}\\
 \class{N/D}_{\Q^\times/\Q^{\times2}}&=\class 3.        \label{eq:ratio3}
\end{align}
In particular, \(N\) is neither a rational square nor in the sign-field
class.

\begin{theorem}[Full Kummer rank]\label{thm:rank-nine}
The nine conjugates \(\beta_1,\ldots,\beta_9\) are linearly independent in
\(L^\times/L^{\times2}\).  Equivalently, \(R=0\).
\end{theorem}

\begin{proof}
Proposition~\ref{prop:one-relation} leaves only \(\ones\).  If
\(\ones\in R\), then
\[
 \prod_{i=1}^{9}\beta_i=\Norm_{K/\Q}(\beta)=N
\]
would be a square in \(L\).  Lemma~\ref{lem:sign-field} and
\eqref{eq:nclass}--\eqref{eq:ratio3} rule this out.  Hence
\(\ones\notin R\), and \(R=0\).
\end{proof}

This proof uses one local row, not a class-group computation or a search
over all \(2^9\) possible products.  The full \(S_9\) symmetry is what
turns a weight-two local observation into a rank-nine global theorem.
